\begin{table}[H]
\centering
\captionsetup{justification=centering}
\caption{Teste de diagnóstico do resíduo para o total de ocupados - modelos univariado e multivariado}
\label{tab:diagocup}
\scalebox{1}{
\renewcommand{\arraystretch}{0.7}
\begin{tabular}{lccc}
\toprule
& \multicolumn{3}{c}{\textbf{Modelo Univariado}} \\
\cmidrule(lr){2-4}
\textbf{Estrato geográfico} & Shapiro-Wilk & Ljung-Box & H \\
\midrule
01 - Belo Horizonte & 0,0000*** & 0,1767 & 0,5274 \\
02 - Entorno e Colar Metropolitano de BH & 0,0037*** & 0,0527* & 0,1645 \\
03 - Sul de Minas & 0,0153** & 0,9479 & 0,9638 \\
04 - Triângulo Mineiro & 0,0751* & 0,4637 & 0,0711* \\
05 - Zona da Mata & 0,7568 & 0,5175 & 0,8997 \\
06 - Norte de Minas & 0,0231** & 0,5524 & 0,5957 \\
07 - Vale do Rio Doce & 0,0175** & 0,6971 & 0,0885* \\
08 - Central & 0,0001*** & 0,6827 & 0,2501 \\
\midrule
& \multicolumn{3}{c}{\textbf{Modelo Multivariado}} \\
\cmidrule(lr){2-4}
\textbf{Estrato geográfico} & Shapiro-Wilk & Ljung-Box & H \\
\midrule
01 - Belo Horizonte & 0,0000*** & 0,1717 & 0,5724 \\
02 - Entorno e Colar Metropolitano de BH & 0,0046*** & 0,0420** & 0,1517 \\
03 - Sul de Minas & 0,0102** & 0,9575 & 0,9656 \\
04 - Triângulo Mineiro & 0,0839* & 0,5055 & 0,0808* \\
05 - Zona da Mata & 0,5744 & 0,5625 & 0,8753 \\
06 - Norte de Minas & 0,0222** & 0,6040 & 0,5003 \\
07 - Vale do Rio Doce & 0,0538* & 0,7628 & 0,1688 \\
08 - Central & 0,0001*** & 0,6369 & 0,2324 \\
\bottomrule
\end{tabular}}
\fonte{Elaboração própria, com base nos dados da PNAD Contínua. Nota: * representa rejeição de \(H_0\) a 10\%, ** a 5\% e *** a 1\%. As oito primeiras observações foram descartadas, em razão do período de estabilização do filtro de Kalman.}
\end{table}
